\expandafter\ifx\csname CAPETINCLUDEMODE\endcsname\relax
\documentclass[11pt,a4paper]{article}
% preamble.tex — packages et commandes communes a toutes les fiches CAPET SII IE
\input{glyphtounicode}
\pdfgentounicode=1
\usepackage{cmap}
\usepackage[T1]{fontenc}
\usepackage[utf8]{inputenc}
\usepackage[french]{babel}
% babel[french] active a \begin{document} des raccourcis de ponctuation (:;!?)
% qui cassent la syntaxe to[...] de circuitikz (erreur "+ or - expected") ;
% on les desactive, ce qui ne change que l'espacement fin avant ces signes.
\AtBeginDocument{\shorthandoff{:;!?}}
\usepackage[a4paper, top=1.6cm, bottom=1.6cm, left=1.4cm, right=1.4cm]{geometry}
\setlength{\headheight}{14.5pt}
\usepackage{amsmath, amssymb, amsfonts}
\usepackage{siunitx}
\sisetup{locale=FR, output-decimal-marker={,}, inter-unit-product=\ensuremath{{}\cdot{}},
         group-minimum-digits=4}
% Unites non fournies par siunitx, declarees ici une fois pour toutes (fiches + exercices).
% Les gardes \ifdefined evitent tout conflit avec une redeclaration locale restee dans un
% fichier : ne pas retirer les gardes, ne pas redeclarer ces unites dans les fichiers.
\DeclareSIUnit{\jour}{j}
\ifdefined\var\else\DeclareSIUnit{\var}{var}\fi   % puissance reactive
\ifdefined\vad\else\DeclareSIUnit{\vad}{VAD}\fi   % puissance deformante
\ifdefined\tr\else\DeclareSIUnit{\tr}{tr}\fi      % tour (tr/min)
\usepackage{array, booktabs, multirow, tabularx}
\usepackage{xcolor}
\usepackage[most, breakable]{tcolorbox}
\usepackage{multicol}
\usepackage{enumitem}
\usepackage{mathtools}
\usepackage{lmodern}
% Refonte visuelle 2026 (direction "moderne coloré") : Lato pour le texte,
% maths sans-serif assorties via newtxsf (charge apres amsmath/amssymb).
\usepackage[default]{lato}
\usepackage{newtxsf}
\usepackage{microtype}
\usepackage{fancyhdr}
\usepackage{tikz}
\usetikzlibrary{positioning,arrows.meta}
\usepackage{circuitikz}
\usepackage{lastpage}

% ── Couleurs par domaine (palette modernisée 2026) ────────────────────
\definecolor{colElec}{RGB}{29, 78, 216}        % Électrotechnique — bleu vif
\definecolor{colElecLight}{RGB}{219, 234, 254} % fond bleu clair
\definecolor{colAuto}{RGB}{21, 128, 61}        % Automatique — vert
\definecolor{colMeca}{RGB}{194, 65, 12}        % Mécanique — orange
\definecolor{colTherm}{RGB}{185, 28, 28}       % Thermique — rouge
\definecolor{colQT}{RGB}{232, 244, 233}        % (hérité, plus utilisé par questiontype)
\definecolor{colJury}{RGB}{255, 248, 225}      % fond jury (crème)
\definecolor{colJuryFrame}{RGB}{202, 138, 4}   % bordure/bandeau jury (ambre)
\definecolor{colGrayLight}{RGB}{245, 245, 247} % fond gris neutre
\definecolor{colSeparator}{RGB}{180, 180, 180} % gris séparateur
\definecolor{colExo}{RGB}{8, 145, 178}          % Banque d'exercices — cyan (distinct des fiches)
% ── Encre et gris de la refonte ───────────────────────────────────────
\definecolor{colInk}{RGB}{23, 32, 42}          % texte courant (quasi-noir)
\definecolor{colSlate}{RGB}{83, 98, 118}       % gris ardoise (méta, en-têtes)
\definecolor{colRule}{RGB}{203, 213, 225}      % filets clairs
\definecolor{colJuryBg}{RGB}{255, 248, 225}    % fond crème jury
\definecolor{colPresence}{RGB}{83, 98, 118}    % bloc 5 (présence dans les sujets)
\AtBeginDocument{\color{colInk}}

% ── Style encadré exercice + corrigé (banque d'exercices) ─────────────
\tcbset{
  exoBox/.style={
    breakable,
    enhanced,
    colback=colExo!8!white,
    colframe=colExo,
    colbacktitle=colExo,
    coltitle=white,
    fonttitle=\bfseries\small,
    left=4pt, right=4pt, top=3pt, bottom=3pt,
    before skip=5pt, after skip=5pt,
  }
}

% ── Bascule enonce / corrige (banque d'exercices) ─────────────────────
% \CAPETCORRIGE est defini en ligne de commande (pdflatex "\def\CAPETCORRIGE{1}\input{...}")
% pour compiler la variante corriges d'un meme fichier source ; par defaut (non defini),
% la compilation produit la variante enonces. \CAPETSHOW{enonce}{corrige} choisit la branche.
%
% TROIS MODES (refonte 2026-09-20, manuel de revision) :
%   - aucun drapeau   -> imprime l'enonce            (variante enonces, INCHANGE)
%   - \CAPETCORRIGE   -> imprime le corrige          (variante corriges, INCHANGE)
%   - \CAPETMANUEL    -> imprime l'enonce ET met le corrige de cote pour l'annexe
%
% Le mode manuel stocke chaque corrige dans une macro (\csgdef), sans ecriture
% fichier ni re-tokenisation : le contenu (circuitikz, accents, %) traverse intact.
% L'annexe les rejoue ensuite dans l'ordre via \CAPETANNEXE.
\usepackage{etoolbox}

\newcounter{CAPETexo}          % numero continu d'exercice dans le manuel
\setcounter{CAPETexo}{0}
\newcounter{CAPETannexe}       % compteur de relecture en annexe

% Titre courant de la notion en cours (pose par \CAPETNOTION), sert d'intertitre annexe
\newcommand{\CAPETnotioncode}{}
\newcommand{\CAPETnotiontitre}{}

% \CAPETNOTION{code}{titre} : ouvre une notion dans le manuel (memorise le libelle,
% et marque le rang du premier exercice de la notion pour l'intertitre d'annexe).
\newcommand{\CAPETNOTION}[2]{%
  \renewcommand{\CAPETnotioncode}{#1}%
  \renewcommand{\CAPETnotiontitre}{#2}%
  \csgdef{CAPETnotionstart\the\numexpr\value{CAPETexo}+1\relax}{#1 -- #2}%
}

% ── \degree robuste ───────────────────────────────────────────────────
% Certains fichiers d'exercices utilisent \degree. Ni le preambule ni siunitx
% ne le definissent, et \providecommand local ne suffit pas : sous hyperref,
% \degree existe deja comme commande factice (\ERROR), donc \providecommand
% s'abstient et l'appel echoue a la compilation du manuel (erreur invisible
% en compilation isolee). On le (re)definit donc inconditionnellement ici,
% apres le chargement des packages, pour tout le corpus.
\AtBeginDocument{\providecommand{\degree}{}\renewcommand{\degree}{\ensuremath{^\circ}}}

% Libelle decoratif : les fichiers d'exercices utilisent aussi \CAPETSHOW dans
% leur \fancyhead et leur bloc de titre, pour afficher "(enonces)" ou "(corriges)".
% Ces appels-la ne sont PAS des exercices : en mode manuel ils ne doivent ni
% consommer un numero, ni alimenter l'annexe. On les reconnait a leur premier
% argument, strictement egal au libelle ci-dessous.
\newcommand{\CAPETLIBELLEENONCES}{(\'enonc\'es)}

\ifdefined\CAPETMANUEL
  % --- Mode manuel : enonce en place, corrige differe vers l'annexe ---
  \newcommand{\CAPETSHOW}[2]{%
    \def\CAPETargun{#1}%
    \ifx\CAPETargun\CAPETLIBELLEENONCES
      #1%            appel decoratif : rendu tel quel, hors numerotation
    \else
      \CAPETSHOWEXO{#1}{#2}%
    \fi
  }
  \newcommand{\CAPETSHOWEXO}[2]{%
    \stepcounter{CAPETexo}%
    \csgdef{CAPETcorrige\theCAPETexo}{#2}%
    \edef\CAPETtmpnum{\theCAPETexo}%
    #1%
    \par\vspace{2pt}%
    {\footnotesize\itshape\color{colExo}%
     \hyperref[capet:cor:\CAPETtmpnum]{Corrige de l'exercice \CAPETtmpnum\ $\rightarrow$ p.~\pageref{capet:cor:\CAPETtmpnum}}%
     \label{capet:eno:\CAPETtmpnum}}%
  }
  % \CAPETANNEXE : rejoue tous les corriges captures, dans l'ordre, avec intertitres
  \newcommand{\CAPETANNEXE}{%
    \setcounter{CAPETannexe}{0}%
    \loop\ifnum\value{CAPETannexe}<\value{CAPETexo}%
      \stepcounter{CAPETannexe}%
      \edef\CAPETcurnum{\theCAPETannexe}%
      \ifcsname CAPETnotionstart\CAPETcurnum\endcsname
        \par\medskip
        {\normalsize\bfseries\color{colExo}\csuse{CAPETnotionstart\CAPETcurnum}}%
        \par\nobreak\smallskip
      \fi
      \begin{tcolorbox}[exoBox, breakable,
        title={Corrige de l'exercice \CAPETcurnum\
               \hfill\normalfont\footnotesize
               \hyperref[capet:eno:\CAPETcurnum]{enonce p.~\pageref{capet:eno:\CAPETcurnum}}}]%
      \label{capet:cor:\CAPETcurnum}%
      \csuse{CAPETcorrige\CAPETcurnum}%
      \end{tcolorbox}%
    \repeat
  }
\else
  % --- Modes historiques, comportement strictement inchange ---
  \newcommand{\CAPETSHOW}[2]{\ifdefined\CAPETCORRIGE #2\else #1\fi}
  \newcommand{\CAPETANNEXE}{}
\fi

% ── Difficulté 1 à 3 (étoiles pleines, amssymb) ───────────────────────
\newcommand{\difficulte}[1]{%
  \ifcase#1\relax
  \or $\bigstar$\;niveau 1/3%
  \or $\bigstar\bigstar$\;niveau 2/3%
  \or $\bigstar\bigstar\bigstar$\;niveau 3/3%
  \fi
}

% ── Style encadré extrait de document technique (banque d'exercices) ──
\tcbset{
  docTech/.style={
    breakable,
    enhanced,
    colback=colGrayLight,
    colframe=colSeparator,
    fonttitle=\bfseries\small\itshape,
    coltitle=black,
    colbacktitle=colGrayLight,
    left=4pt, right=4pt, top=3pt, bottom=3pt,
    before skip=4pt, after skip=4pt,
  }
}

% ── Style tcolorbox de section principale ────────────────────────────
% Refonte 2026 : bandeau titre coloré arrondi + fond très pâle, sans cadre.
\tcbset{
  ficheSection/.style={
    breakable, enhanced jigsaw,
    frame hidden, boxrule=0pt, arc=1.8mm,
    colback=#1!4!white,
    colbacktitle=#1,
    coltitle=white,
    fonttitle=\sffamily\bfseries,
    left=7pt, right=7pt, top=5pt, bottom=5pt,
    toptitle=3pt, bottomtitle=3pt,
    before skip=10pt, after skip=8pt,
  }
}

% ── Style encadré question type ──────────────────────────────────────
% Carte blanche, bandeau sombre, liseré très clair.
\tcbset{
  questiontype/.style={
    breakable, enhanced jigsaw,
    colback=white,
    frame hidden, boxrule=0pt, arc=1.8mm,
    colbacktitle=colElec!78!black,
    coltitle=white,
    fonttitle=\sffamily\bfseries\small,
    left=7pt, right=7pt, top=4pt, bottom=4pt,
    toptitle=2.5pt, bottomtitle=2.5pt,
    before skip=8pt, after skip=8pt,
    borderline={0.6pt}{0pt}{colElec!12!white},
  }
}

% ── Style encadré jury ───────────────────────────────────────────────
% Bandeau ambre, fond crème, sans cadre.
\tcbset{
  jury/.style={
    breakable, enhanced jigsaw,
    colback=colJuryBg,
    frame hidden, boxrule=0pt, arc=1.8mm,
    colbacktitle=colJuryFrame,
    coltitle=white,
    fonttitle=\sffamily\bfseries,
    left=7pt, right=7pt, top=4pt, bottom=4pt,
    toptitle=3pt, bottomtitle=3pt,
    before skip=10pt, after skip=8pt,
  }
}

% ── En-tête / pied de page communs (refonte 2026) ─────────────────────
% #1 code fiche, #2 titre court, #3 couleur domaine
\newcommand{\ficheHeaders}[3]{%
  \fancyhead[L]{\small\sffamily{\bfseries\textcolor{#3}{#1}}\hspace{0.7em}\textcolor{colSlate}{#2}}%
  \fancyhead[R]{\small\sffamily\textcolor{colSlate}{CAPET SII IE\, --- \,Session 2026}}%
  \fancyfoot[C]{\small\sffamily\textcolor{colSlate}{\thepage\,/\,\pageref{LastPage}}}%
  \fancyhead[C]{}\fancyfoot[L]{}\fancyfoot[R]{}%
}

% ── Bloc en-tête de fiche (BLOC 0, refonte 2026) ──────────────────────
% #1 code  #2 titre  #3 couleur domaine  #4 domaine  #5 ligne présence  #6 sous-thèmes
\newcommand{\ficheEntete}[6]{%
  \begin{tcolorbox}[enhanced, breakable,
    frame hidden, boxrule=0pt, arc=2.5mm,
    interior style={left color=#3, right color=#3!75!black},
    left=9pt, right=9pt, top=7pt, bottom=7pt,
    before skip=0pt, after skip=6pt]
    {\color{white}\sffamily\bfseries\LARGE #1\quad #2}\\[4pt]
    {\color{#3!22!white}\sffamily\small \MakeUppercase{#4}\quad\textbar\quad #5}\\[1pt]
    {\color{#3!35!white}\sffamily\footnotesize Sous-thèmes : #6}
  \end{tcolorbox}%
}

% ── Mise en page ─────────────────────────────────────────────────────
\pagestyle{fancy}
\fancyhf{}
\renewcommand{\headrulewidth}{0.5pt}
\renewcommand{\footrulewidth}{0pt}
\AtBeginDocument{%
  \renewcommand{\headrule}{{\color{colRule}\hrule width\headwidth height 0.5pt}}%
}

\setlist[itemize]{leftmargin=1.2em, itemsep=1pt, topsep=2pt, parsep=0pt}
\setlist[enumerate]{leftmargin=1.2em, itemsep=3pt, topsep=2pt, parsep=0pt}

\setlength{\parskip}{2pt}
\setlength{\columnsep}{10pt}

\def\CAPETBEGINFILE{\begin{document}}
\def\CAPETENDFILE{\end{document}}
\else
\def\CAPETBEGINFILE{}
\def\CAPETENDFILE{}
\fi

\ficheHeaders{T14}{Variation de vitesse}{colElec}

\CAPETBEGINFILE

% -----------------------------------------------------------------------
%  BLOC EN-TETE
% -----------------------------------------------------------------------

\ficheEntete{T14}{Variation de vitesse et de couple des machines}{colElec}%
  {électrotechnique - Variateurs}%
  {\textbf{Présence :} \textbf{5/13} sujets (2013-2026) - thème \emph{régulier}}%
  {commande scalaire V/f - THD et commande MLI - variateur de fréquence ATV630 - boucle de vitesse PI - freinage régénératif 4 quadrants - câblage variateur-automate}

% -----------------------------------------------------------------------
%  1. RAPPELS DE COURS
% -----------------------------------------------------------------------

\begin{tcolorbox}[
  ficheSection=colElec,
  title={1. Rappels de cours},
]

% -- 1.1 Variateur de frequence ----------------------------------------
\textbf{1.1 Structure d'un variateur de fréquence}

\smallskip
\begin{center}
\begin{tikzpicture}[>=Stealth, font=\small, node distance=8mm]
  \node[draw, align=center, minimum width=20mm, minimum height=10mm] (redr) {Redresseur};
  \node[draw, align=center, minimum width=20mm, minimum height=10mm, right=13mm of redr] (filt) {Bus DC\\filtré};
  \node[draw, align=center, minimum width=20mm, minimum height=10mm, right=13mm of filt] (ond) {Onduleur\\MLI};
  \node[left=13mm of redr] (res) {Réseau AC};
  \node[right=13mm of ond] (mas) {MAS};
  \draw[->] (res) -- (redr);
  \draw[->] (redr) -- (filt);
  \draw[->] (filt) -- (ond);
  \draw[->] (ond) -- (mas);
\end{tikzpicture}
\end{center}
Avantages : contrôle de la vitesse et du couple, démarrage progressif (rampe d'accélération), freinage par injection de courant continu ou récupératif, économie d'énergie sur charges à couple variable ($P \propto N^3$).

\medskip\hrule\vspace{4pt}

% -- 1.2 Commande scalaire V/f ----------------------------------------
\textbf{1.2 Commande scalaire V/f}

\smallskip
\[
  \frac{V_s}{f_s} = \text{constante} \quad \Rightarrow \quad \text{flux magnétique constant}
\]
Principe : faire varier simultanément la fréquence statorique $f_s$ et la tension $V_s$ pour maintenir le flux et donc le couple disponible.
Boucle de vitesse avec correcteur PI :
\[
  H_{BF}(p) = \frac{1}{1 + \dfrac{2m}{\omega_0}\,p + \dfrac{1}{\omega_0^2}\,p^2}
  \qquad;\qquad
  m = \text{amortissement}, \quad \omega_0 = \text{pulsation naturelle}
\]
Données 2024 : $K_p = 0{,}35$, $K_i = 2{,}5$, erreur statique nulle, dépassement nul.

\medskip\hrule\vspace{4pt}

% -- 1.3 THD et choix de commande ------------------------------------
\textbf{1.3 Taux de distorsion harmonique (THD) - commande pleine onde vs MLI}

\smallskip
\[
  THD = \frac{\sqrt{I_2^2 + I_3^2 + \cdots + I_n^2}}{I_1}
  \qquad;\qquad
  \text{Norme IEC/EN 61000-2-2 : } THD < 25\,\%
\]

\begin{center}
\renewcommand{\arraystretch}{1.25}
\begin{tabular}{lll}
\toprule
\textbf{Commande} & \textbf{THD typique} & \textbf{Conformité norme} \\
\midrule
Pleine onde & $\approx 61\,\%$ & Non (> 25 \%) \\
MLI ($f_p = \SI{6}{\kilo\hertz}$) & $\approx 25\,\%$ & Limite acceptable \\
\bottomrule
\end{tabular}
\end{center}
La MLI repousse les harmoniques vers les hautes fréquences (autour de $f_p$ et ses multiples) : THD du fondamental amélioré.

\medskip\hrule\vspace{4pt}

% -- 1.4 Dimensionnement variateur ------------------------------------
\textbf{1.4 Dimensionnement et câblage du variateur ATV630}

\smallskip
Courant nominal moteur : $I_n = \dfrac{P_u}{\eta\,\sqrt{3}\,U_{comp}\,\cos\varphi}$

\begin{center}
\renewcommand{\arraystretch}{1.25}
\begin{tabular}{lll}
\toprule
\textbf{Critère de choix} & \textbf{Paramètre} & \textbf{Exemple ATV630C11N4} \\
\midrule
Puissance moteur & $P_u$ & 90\,kW / 100\,kW \\
Surcharge admissible & ND (110\,\%) ou HD (150\,\%) & HD si charge lourde \\
Tension réseau & 400\,V triphasé & - \\
Pouvoir de coupure disjoncteur amont & $I_{cu}$ & 50\,kA (NSX250N MA220) \\
\bottomrule
\end{tabular}
\end{center}
\textbf{Câblage :} sens de rotation (LI1 / LI2), référence de vitesse (entrée 0-10\,V ou 4-20\,mA), retour de défaut à sortie DO. Rapport jury 2021 : maîtrise du câblage automate-variateur obligatoire.

\end{tcolorbox}

% -----------------------------------------------------------------------
%  2. FORMULAIRE SYNTHETIQUE
% -----------------------------------------------------------------------

\begin{tcolorbox}[
  ficheSection=colElec,
  title={2. Formulaire synthétique},
]
\small
\begin{multicols}{2}
\renewcommand{\arraystretch}{1.45}
\begin{tabular}{@{}p{4.5cm}ll@{}}
\toprule
\textbf{Formule} & \textbf{Usage} & \textbf{U.} \\
\midrule
\multicolumn{3}{@{}l}{\footnotesize\textit{Commande V/f}} \\
$V_s/f_s = \mathrm{cste}$ & Loi scalaire & - \\
$N_s = 60f/p$ & Synchro & tr/min \\
\midrule
\multicolumn{3}{@{}l}{\footnotesize\textit{THD}} \\
$THD = \sqrt{\sum_{k\geq2} I_k^2}/I_1$ & Distorsion & - \\
$THD < 25\,\%$ & Norme & - \\
\midrule
\multicolumn{3}{@{}l}{\footnotesize\textit{Courant moteur}} \\
$I_n = P_u/(\eta\sqrt{3}U\cos\varphi)$ & Nominal & A \\
\bottomrule
\end{tabular}

\columnbreak

\renewcommand{\arraystretch}{1.45}
\begin{tabular}{@{}p{4.5cm}ll@{}}
\toprule
\textbf{Formule} & \textbf{Usage} & \textbf{U.} \\
\midrule
\multicolumn{3}{@{}l}{\footnotesize\textit{Boucle vitesse PI}} \\
$H(p) = K_p(1 + 1/(\tau_i p))$ & Correcteur & - \\
$m$ = amortissement & BF 2\textsuperscript{e} ordre & - \\
\midrule
\multicolumn{3}{@{}l}{\footnotesize\textit{Quadrants}} \\
I : $C>0,\Omega>0$ & Moteur direct & - \\
II : $C<0,\Omega>0$ & Freinage récup. & - \\
III,IV & Rotation inverse & - \\
\midrule
\multicolumn{3}{@{}l}{\footnotesize\textit{Données 2024}} \\
$f_p = \SI{6}{\kilo\hertz}$ & MLI & Hz \\
$THD_{pleine} \approx 61{,}4\,\%$ & Pleine onde & - \\
\bottomrule
\end{tabular}
\end{multicols}
\normalsize
\end{tcolorbox}

% -----------------------------------------------------------------------
%  3. QUESTIONS TYPES
% -----------------------------------------------------------------------

\begin{tcolorbox}[
  ficheSection=colElec,
  title={3. Questions types - extraites des sujets réels},
]

% -- QT1 -------------------------------------------------------------

\begin{tcolorbox}[questiontype,
  title={QT1 - THD et conformité norme - commande MLI vs pleine onde \hfill\normalfont\footnotesize E1 2024 (B10-B11)}
]
\textbf{Énoncé.}
Le DT3 donne le spectre du courant réseau en pleine onde : $I_1 = 170\,\mathrm{A}$, $I_3 = 95\,\mathrm{A}$, $I_5 = 27\,\mathrm{A}$, $I_7 = 27\,\mathrm{A}$, $I_9 = 20\,\mathrm{A}$. En MLI ($f_p = 6\,\mathrm{kHz}$) : $I_1 = 170\,\mathrm{A}$, $I_3 = 37\,\mathrm{A}$, $I_5 = 22\,\mathrm{A}$. Calculer le THD dans chaque cas et conclure.

\tcbline

\textbf{Corrigé.}

\textbf{Étape 1 — Pleine onde.}
\[
  THD = \frac{\sqrt{95^2 + 27^2 + 27^2 + 20^2}}{170} = \frac{\sqrt{9\,025 + 729 + 729 + 400}}{170} = \frac{\sqrt{10\,883}}{170} \approx \mathbf{61{,}4\,\%}
\]
Non conforme (> 25\,\%).

\textbf{Étape 2 — MLI.}
\[
  THD = \frac{\sqrt{37^2 + 22^2}}{170} = \frac{\sqrt{1\,369 + 484}}{170} = \frac{\sqrt{1\,853}}{170} \approx \mathbf{25{,}3\,\%}
\]

\textbf{Conclusion.} La commande MLI est légèrement \textbf{non conforme} au seuil strict de $25\,\%$ (écart $\approx 0{,}3$ point) ; une conclusion de conformité ne serait acceptable que si une tolérance de mesure est explicitement justifiée. La MLI reste néanmoins indispensable pour réduire fortement la distorsion par rapport à la pleine onde.
\end{tcolorbox}

\vspace{4pt}

% -- QT2 -------------------------------------------------------------

\begin{tcolorbox}[questiontype,
  title={QT2 - Choix et dimensionnement du variateur ATV630 \hfill\normalfont\footnotesize Rapport jury 2021 (Q15-Q18)}
]
\textbf{Énoncé.}
Choisir le modèle ATV630 pour un moteur de $P_u = 90\,\mathrm{kW}$, $U = 400\,\mathrm{V}$, $\cos\varphi = 0{,}86$, $\eta = 0{,}942$, surcharge 150\,\% (mode HD). Indiquer le disjoncteur amont ($I_{cu,reseau} = 50\,\mathrm{kA}$).

\tcbline

\textbf{Corrigé.}

\textbf{Étape 1 — Courant nominal moteur.}
\[
  I_n = \frac{P_u}{\eta\,\sqrt{3}\,U\,\cos\varphi} = \frac{90\,000}{0{,}942 \times 1{,}732 \times 400 \times 0{,}86} \approx \frac{90\,000}{562} \approx 160\,\mathrm{A}
\]

\textbf{Étape 2 — Choix variateur.} ATV630\textbf{C11N4} ($I_n = 170\,\mathrm{A}$, 110\,kW, HD 150\,\%, 400\,V). La puissance 110\,kW > 90\,kW et le courant 170\,A > 160\,A.

\textbf{Étape 3 — Disjoncteur amont.} NSX250N MA220 - déclencheur \textbf{MA} (magnétique seul, seuil $I_{rm} = \SI{2860}{\ampere}$, le variateur assure la protection thermique), $I_{cu} = 50\,\mathrm{kA}$ (adapté au $I_{cc}$ réseau).

\textbf{Conclusion.} Le couple variateur + protection amont respecte simultanément les trois critères du cahier des charges (courant nominal, surcharge HD, pouvoir de coupure) : le choix est validé.
\end{tcolorbox}

\vspace{4pt}

% -- QT3 -------------------------------------------------------------

\begin{tcolorbox}[questiontype,
  title={QT3 - Boucle de vitesse scalaire - correcteur PI \hfill\normalfont\footnotesize E1 2024 (C5-C11)}
]
\textbf{Énoncé.}
La boucle de vitesse de la MAS (commande scalaire) associe, en retour unitaire, un correcteur PI $C(p) = K_p + \dfrac{K_i}{p}$ et le processus mécanique $G(p) = \dfrac{K_{meca}}{1 + \tau_{meca}\,p}$ avec $K_{meca} = 250$ et $\tau_{meca} = 8{,}5\,\si{\second}$. Déterminer la fonction de transfert en boucle fermée $H_{BF}(p)$, puis exprimer $\omega_0$ et $m$ en fonction de $K_p$, $K_i$, $K_{meca}$ et $\tau_{meca}$. Conclure avec les valeurs retenues au corrigé : $K_p = 0{,}35$, $K_i = 2{,}5$.

\tcbline

\textbf{Corrigé.} \textit{[corrigé reconstitué : mise en équation standard, le détail du sujet n'étant pas disponible]}

\textbf{Étape 1 — Boucle fermée.} La chaîne directe vaut $C(p)\,G(p) = \dfrac{K_{meca}(K_p\,p + K_i)}{p\,(1 + \tau_{meca}\,p)}$, d'où :
\[
  H_{BF}(p) = \frac{C(p)\,G(p)}{1 + C(p)\,G(p)}
  = \frac{K_{meca}(K_p\,p + K_i)}{\tau_{meca}\,p^2 + (1 + K_{meca}K_p)\,p + K_{meca}K_i}
\]
Le dénominateur, divisé par $K_{meca}K_i$, s'identifie à la forme canonique $1 + \dfrac{2m}{\omega_0}\,p + \dfrac{1}{\omega_0^2}\,p^2$.

\textbf{Étape 2 — Identification.}
\[
  \omega_0 = \sqrt{\frac{K_{meca}\,K_i}{\tau_{meca}}}
  \qquad;\qquad
  m = \frac{1 + K_{meca}K_p}{2\,\sqrt{K_{meca}\,K_i\,\tau_{meca}}}
\]
Application numérique avec $K_p = 0{,}35$ et $K_i = 2{,}5$ :
\[
  \omega_0 = \sqrt{\frac{250 \times 2{,}5}{8{,}5}} \approx \SI{8,6}{\radian\per\second}
  \qquad;\qquad
  m = \frac{1 + 87{,}5}{2\,\sqrt{250 \times 2{,}5 \times 8{,}5}} \approx 0{,}61
\]

\textbf{Conclusion.} L'intégrateur du PI garantit une erreur statique nulle en régime permanent, et le corrigé officiel retient $K_p \approx 0{,}35$, $K_i \approx 2{,}5$ en annonçant un dépassement nul : le cahier des charges (erreur statique nulle, dépassement nul) est ainsi satisfait par un unique correcteur.
\end{tcolorbox}

\end{tcolorbox}

% -----------------------------------------------------------------------
%  4. ATTENTES DU JURY
% -----------------------------------------------------------------------

\begin{tcolorbox}[
  jury,
  title={4. Attentes du jury - Erreurs récurrentes et points de vigilance},
]

\begin{enumerate}[label=\textbf{\arabic*.}, leftmargin=*, itemsep=5pt]

  \item \textbf{Justifier le choix du variateur - pas seulement la référence.}
  Citer les critères : puissance ($P_u$), type de service (ND ou HD), tension réseau, pouvoir de coupure. Rapport jury 2021 : "sans justification, sans indiquer clairement les critères de choix".

  \item \textbf{Disjoncteur MA devant un variateur.}
  Le variateur assure lui-même la protection thermique du moteur. Utiliser un déclencheur magnétique seul (MA) en amont. Un disjoncteur magnéto-thermique crée une double protection thermique - décoordination.

  \item \textbf{THD : mettre au carré sous la racine.}
  Oublier de mettre les harmoniques au carré, ou diviser par l'amplitude totale au lieu de $I_1$, sont les erreurs les plus fréquentes.

  \item \textbf{Commande scalaire vs vectorielle.}
  La commande scalaire V/f suffit pour les charges à couple variable (pompe, ventilateur). La commande vectorielle FOC n'est pas exigée au CAPET mais peut être citée pour les applications à dynamique élevée.

  \item \textbf{Câblage automate-variateur.}
  Rapport jury 2021 : "un candidat au CAPET doit être capable de câbler une entrée ou une sortie TOR-analogique d'un automate industriel programmable, ainsi que la partie commande d'un variateur de fréquence".

  \item \textbf{Freinage régénératif et 4 quadrants.}
  Identifier les quadrants (couple/vitesse). Un onduleur 4 quadrants est nécessaire pour le freinage récupératif (quadrant II ou IV). Vérifier que la structure matérielle le permet.

\end{enumerate}

\medskip
\begin{tcolorbox}[
  breakable,
  colback=colJuryFrame!12!colJury,
  colframe=colJuryFrame,
  left=4pt, right=4pt, top=2pt, bottom=2pt,
  fontupper=\small\itshape
]
Rapport jury 2021 : «~Le jury constate que beaucoup de candidats ne sont pas à l'aise dans la réalisation de schémas électriques. Un candidat au CAPET doit être capable de câbler une entrée ou une sortie TOR-analogique d'un automate industriel programmable, ainsi que la partie commande d'un variateur de fréquence.~»\\
Rapport jury 2024 : «~Le contrôle de la vitesse peut être réalisé via une commande scalaire de la machine asynchrone et une boucle d'asservissement. L'objectif est d'étudier les paramètres à respecter : erreur statique, dépassement et marge de phase.~»
\end{tcolorbox}

\end{tcolorbox}

% -----------------------------------------------------------------------
%  PIED - PRESENCE DANS LES SUJETS
% -----------------------------------------------------------------------

\begin{tcolorbox}[
  ficheSection=colPresence,
  title={Présence dans les sujets analysés (2013-2026)},
]
\begin{center}
\footnotesize\renewcommand{\arraystretch}{1.15}
\begin{tabular}{lccccccccccccc}
\toprule
\textbf{Année} & 13 & 14 & 15 & 16 & 17 & 18 & 19 & 20 & 21 & 22 & 23 & 24 & 25 \\
\midrule
E1 (commune)   & - & - & - & - & - & - & - & $\checkmark$ & - & - & - & $\checkmark$ & - \\
E2 discipl. IE & - & - & - & $\checkmark$ & - & $\checkmark$ & - & - & $\checkmark$ & - & - & - & - \\
E2 appl. comm. & - & - & - & - & - & - & - & - & - & - & - & - & - \\
\bottomrule
\end{tabular}
\end{center}
\smallskip
\footnotesize\textbf{Sous-thèmes fréquents :}
commande V/f et THD (4/13) - choix/câblage variateur ATV (3/13) - boucle de vitesse PI (2/13) - freinage 4 quadrants (sujets 2026) - dimensionnement disjoncteur MA (3/13)

\medskip
\textbf{Ordres de grandeur (repères rapides).}
\begin{itemize}[leftmargin=1.3em]
  \item Réseau BT FR : tension phase-neutre typique \SI{230}{\volt}, tension phase-phase typique \SI{400}{\volt}, fréquence \SI{50}{\hertz}.
  \item Courants d'étude fréquents en sujet : de quelques \si{\ampere} à quelques dizaines de \si{\ampere}.
  \item En conversion d'énergie, un rendement attendu reste généralement entre 0{,}8 et 0{,}98 selon la technologie.
\end{itemize}

\medskip
\textbf{Synthèse express (à refaire sans support).}
\begin{enumerate}[leftmargin=1.4em]
  \item Poser les hypothèses du modèle et les conventions de signe.
  \item Écrire 2 à 4 relations clés avec unités.
  \item Réaliser un calcul numérique complet et vérifier l'ordre de grandeur.
  \item Justifier le résultat par un argument technique de niveau CAPET.
\end{enumerate}

\end{tcolorbox}

\CAPETENDFILE
